\begin{abstract}
Deep learning models for visual recognition often exhibit systematic errors due to underrepresented semantic subpopulations. While existing debugging frameworks can identify these failure slices, effectively repairing them remains difficult. Current solutions often rely on manually designed prompts to generate synthetic images—an approach that introduces distribution shift and semantic errors, often creating new bugs. To address these issues, we introduce \model, a framework for distribution-consistent model repair via controlled counterfactual generation. Our approach employs a structural-conditioned diffusion model to generate semantically faithful images that modify only specific failure attributes while preserving the underlying data distribution. To ensure the reliability of the repair data, we implement a closed-loop verification mechanism using a large vision--language model (LVLM) to enforce semantic consistency and label preservation. By retraining models on this curated synthetic data, we significantly reduce errors in rare cases without sacrificing overall performance. Our experiments show that \model achieves superior robustness compared to existing augmentation methods by maintaining high precision in attribute editing.
\end{abstract}

\section{Introduction}
Despite achieving high accuracy on standard benchmarks, computer vision models often fail on rare semantic subpopulations~\cite{barbu2019objectnet,gao2023adaptive,leclerc20223db}. These failures usually stem from dataset bias, where specific attribute combinations (e.g., darker-skinned individuals with specific hair textures) are underrepresented in the training set~\cite{buolamwini2018gender}. Identifying and repairing these "bug slices" is essential for the reliable deployment of AI systems in safety-critical applications.

Recent research has focused on identifying these failures using interpretable debugging pipelines~\cite{gao2023adaptive,chen2023hibug,singla2024data}. For instance, \hibug~\cite{chen2023hibug} uses vision--language models to discover failure slices based on shared visual attributes. However, the subsequent repair process remains a major bottleneck. Existing repair strategies generally fall into two categories: retrieval-based and generation-based. Retrieval-based methods~\cite{gao2023adaptive,singla2024data} collect samples from external datasets, which often introduces a domain gap. Generation-based methods, such as those suggested by \hibug, use text-to-image models to synthesize training data. However, standard generative augmentation suffers from two critical flaws: (1) \textbf{Semantic Drift}: Generative models often fail to preserve attributes unrelated to the intended edit, and (2) \textbf{Distribution Mismatch}: Synthetic images may not align with the original data distribution, leading the model to learn "generative artifacts" rather than the intended semantic concepts.

To overcome these challenges, we propose \model, a targeted model repair framework that treats repair as a \textbf{controlled counterfactual generation} problem. Unlike prior works that use generative models as black-box data sources, \model integrates spatial constraints and semantic feedback to ensure the safety and precision of the repair process. The core of our method lies in two technical insights. First, we use a structural-conditioned diffusion mechanism (e.g., ControlNet) to anchor the generation on existing training instances. This ensures that only the problematic attributes are modified while the background and non-targeted features remain unchanged, preserving distributional consistency. Second, we introduce a semantic guardrail using Large Vision-Language Models (LVLMs) to verify the generated images. This closed-loop system filters out samples where the generative model fails to render the requested attribute or unintentionally alters the class label.

The contributions of this work are as follows:
\begin{itemize}
\item We formulate a targeted model repair pipeline, \model, which moves beyond simple data augmentation by framing the problem as distribution-consistent counterfactual generation.
\item We introduce a closed-loop verification mechanism that leverages LVLMs to mitigate the inherent unreliability of generative models, ensuring that the synthesized repair data is both semantically accurate and label-preserving.
\item We demonstrate that our approach significantly enhances model performance on rare-case failure slices across multiple architectures and datasets, outperforming current baselines in both repair effectiveness and the prevention of new bugs.
\end{itemize}